% ************************** Thesis Abstract *****************************
% Use `abstract' as an option in the document class to print only the titlepage and the abstract.
\begin{abstract}

To be filled

% The world's highest energy proton-proton colliding accelerator, the Large Hadron Collider (LHC) experiment at CERN, was proposed and started in 2008 to extend the frontiers of particle physics. The ATLAS experiment at LHC is designed to record the particle collisions at the LHC and search for Higgs bosons and physics beyond the Standard Model. Due to the need to record a large amount of collision data, the trigger system is applied. The trigger system uses the information provided by the readout system for object recognition, such as position and energy. So that only the events that pass a set of criteria are selected. During the Run-2 period (the year 2015-2018), the LHC collided beams at center-of-mass energy $\sqrt{s}=13$ TeV, and ATLAS recorded about 139 fb$^{-1}$ of proton-proton collision data. Since Run-2 ended in 2018, the LHC has been upgraded and will start Run-3 in April 2022. In order to accommodate the update, the ATLAS trigger system requires higher efficiency and better resolution. Hence the \phase upgrade for the ALTAS trigger system is performed.

% The ATLAS Liquid Argon Calorimeter \phase upgrade provides the new trigger system to obtain better energy resolution and efficiency for selecting electrons, photons, $\tau$ leptons, jets, and missing transverse energy. The new readout system applies over 34000 supercells with 10 times higher granularity than the old readout cells and uses optical fibers to transmit the digitized data from the front-end to the back-end. The supercells with multi-layer structure cover the region of $|\eta|<4.9$. The signals read out from the supercells are digitized by the LTDB at the front-end and transmitted to the back-end via optical fibers, where they are calculated by the LATOME board equipped with a high-speed FPGA for energy and timing reconstruction. The computation result is then sent to FEX for object recognition through optical fibers.

% In this thesis, several commissioning works to validate the new trigger system are presented. The first work is connectivity check. It's performed to check whether the supercells are connected correctly from the detector to the back-end. The SSW scan is used to check the mapping of all the 34048 supercells, and the corresponding analysis framework for the SSW scan is developed in this research. Also, by utilizing such an analysis tool, some other channel issues like bad links or bad ADC configuration are also investigated.

% Besides checking the connectivity of supercells, the system latency of the new trigger system is also calibrated in the study. The system latency is required to be well-tuned and fixed for a stable system. The energy and timing computation values cannot be obtained correctly without the calibrated system latency. In the ATLAS experiment, the bunch crossing identification~(BCID) is used for adjusting the system latency. Such work of adjusting the BCID of over 4000 fibers on the LATOME board is done by aligning the BCID of the new trigger system to the BCID of the reference main readout system. And the result of BCID calibration is given at last.

% The research also utilizes the pilot run data taken in October 2021 to perform a full-powered new readout system test with real particles. In the pilot run, the detector managed to take both beam splash and collision data. With these data, the comparison of BCID and the reconstructed \et between the two readouts is made, where \et is computed by the optimal filtering algorithm with the calibrated optimal filtering coefficient. And through this comparison, the thesis gives the evaluation of the commissioning works and the conclusion on whether the new trigger system can achieve the expected performance.



% first letter Z is for Acronyms
% \nomenclature[z-LAr]{LAr}{Liquid Argon}  
% \nomenclature[z-LHC]{LHC}{Large Hadron Collider}
% \nomenclature[z-LS2]{LS2}{the second LHC long Shutdown}
% \nomenclature[z-ATLAS]{ATLAS}{A Toroidal LHC Apparatus}
% \nomenclature[z-BCID]{BCID}{Bunch Crossing Identification}
% \nomenclature[z-ET]{$E_\text{T}^\text{miss}$}{Missing Transverse Momentum}


\end{abstract}
